%----------------------------------------------------------------------------------------
%	PACKAGES AND OTHER DOCUMENT CONFIGURATIONS
%----------------------------------------------------------------------------------------

\documentclass[10pt,stdletter,dateno,sigleft]{newlfm} % Extra options: 'sigleft' for a left-aligned signature, 'stdletternofrom' to remove the from address, 'letterpaper' for US letter paper - consult the newlfm class manual for more options

%\usepackage{charter} % Use the Charter font for the document text
\usepackage{graphics}

\usepackage{standalone}

\newsavebox{\Luiuc}\sbox{\Luiuc}{\parbox[b]{1.75in}{\vspace{0.5in}
\includegraphics[width=1.2\linewidth]{/Users/hectorbahamonde/RU/Job_Market/logo.jpg}}} % Company/institution logo at the top left of the page
\makeletterhead{Uiuc}{\Lheader{\usebox{\Luiuc}}}

\newlfmP{sigsize=50pt} % Slightly decrease the height of the signature field
\newlfmP{addrfromphone} % Print a phone number under the sender's address
\newlfmP{addrfromemail} % Print an email address under the sender's address
\PhrPhone{\texttt{p}} % Customize the "Telephone" text
\PhrEmail{\texttt{e}} % Customize the "E-mail" text

\lthUiuc % Print the company/institution logo

\usepackage{color} % for colors
\usepackage{hyperref} % for links

\hypersetup{
  colorlinks = true,
  urlcolor = blue,
  pdfpagemode = UseNone
}

%----------------------------------------------------------------------------------------
%	YOUR NAME AND CONTACT INFORMATION
%----------------------------------------------------------------------------------------

%\namefrom{\vspace{-3cm}Hector Bahamonde, PhD} % Name

\addrfrom{
{\vspace{-2cm}}\\ % Date
[12pt]
{\bf Hector Bahamonde, PhD}\\
\color[HTML]{003300}{{\bf Postdoctoral Fellow}}\\
\emph{\color[HTML]{003300}{Center For Inter-American Policy \& Research}}\\
\color[HTML]{003300}{Tulane University}\\
205 Richardson, New Orleans, LA 70118
}

\phonefrom{(732) 318-9650} % Phone number

\emailfrom{\href{mailto:bahamondeh@gmail.com}{BahamondeH@gmail.com} \\ 
\texttt{w}: \href{http://www.hectorbahamonde.com}{www.HectorBahamonde.com}\\
\today, download latest version \href{http://github.com/hbahamonde/Job_Market/raw/master/Bahamonde_Cover_Letter.pdf}{\texttt{{\color{red}here}}}}% Link to last version


%----------------------------------------------------------------------------------------
%	ADDRESSEE AND GREETING/CLOSING
%----------------------------------------------------------------------------------------

%\greetto{\vspace{-1cm}Dear Members of the Search Committee,} % Greeting text
%\closeline{\vspace{-1cm}\includegraphics[width=3cm]{/Users/hectorbahamonde/RU/Administracion/signature.pdf} \hspace{-4cm}} % Closing text




%----------------------------------------------------------------------------------------

\begin{document}
\begin{newlfm}

%----------------------------------------------------------------------------------------
%	LETTER CONTENT
%----------------------------------------------------------------------------------------



\vspace{-2cm}{\bf \huge Research Statement}\\

%\vspace{-1cm}
Within comparative politics, my research explores the intersection between inequality and political development, from a political economy perspective. My {\bf job market paper} (\emph{\input{/Users/hectorbahamonde/RU/Dissertation/Papers/Earthquake_Paper/status.txt}\unskip}) is entitled ``\href{https://github.com/hbahamonde/Earthquake_Paper/raw/master/Bahamonde_Earthquake_Paper.pdf}{\input{/Users/hectorbahamonde/RU/Dissertation/Papers/Earthquake_Paper/title.txt}\unskip}.'' In the paper \input{/Users/hectorbahamonde/RU/Dissertation/Papers/Earthquake_Paper/abstract.txt}\unskip.



% main argument of the book
The job market paper is embedded in a {\bf larger research agenda and book manuscript} entitled \emph{``Structural Transformations in Latin America: State Building and Elite Competition 1850-2010.''} The {\bf main argument of the book} is that the economic structural transformation in Latin America, i.e. the secular decline of agriculture and substantial expansion of manufacturing, imposed tight constraints on the way politics was run by the incumbent landowning class. This was a major change due to the advantage the landed elites enjoyed since colonial times. Where the expansion of the industrial sector was weak, post-colonial norms persisted due to institutional inertia, perpetuating the advantaged position of the agricultural class. Leveraging sectoral outputs dating back to 1900 until 2010 for a sample of Latin American countries, I use panel data {\bf methods} (particularly, Cox-proportional hazard models), time series analyses (VAR models, impulse response functions and Granger-causality tests) and fine-grained qualitative data to support my argument. {\bf Another mayor contribution of the book} is measurement. One of the biggest gaps in the literature is the lack of a measurement of state capacities able to capture variations of \emph{stateness} overtime. Using a novel dataset, the book proposes measuring state capacities using earthquake data. The rationale is very intuitive: the enforcement of quake-sensitive building codes throughout the territory is a \emph{reflection} of the overall (\emph{in})capacity of states of solving both logistic and political limitations at the subnational level. \emph{Why does a 7.0 earthquake flatten Haiti leaving at least 100,000 deaths while a 8.8 earthquake in Chile in the same year leaves just 525 deaths?} By exploiting  variations on earthquake death-tolls and local population to weight the number of deaths since 1900, I measure state capacities overtime. Keeping magnitudes constant, casualty differentials should be attributed to the \emph{lack of state capacities}.


In one of the {\bf sections of the book}, I explain that the structural transformation required both sectors to grow in a \emph{balanced} fashion, \emph{leveling both elites in their relative political, economic and military capacities}. Leveraging the dual sector model of economic growth, I sketch a theory of political and economic development that stresses the structural economic dependence of both sectors. The agricultural sector supplies labor and cheap foodstuff which the industrial sector demanded, promoting balanced economic development of the two sectors. The political consequence of balanced economic growth was the mutual political dependence and the need of inter-elite compromises that fostered both economic and political development in the long-run. The paper version of the chapter (currently \emph{\input{/Users/hectorbahamonde/RU/Dissertation/Papers/NegativeLink/status.txt}\unskip}) can be downloaded \href{https://github.com/hbahamonde/Negative_Link_Paper/raw/master/Bahamonde_NegativeLink.pdf}{\texttt{here}}. In another section of the book, I explain how the emergence of a strong industrial sector accelerated the implementation of fiscal institutions. From a fiscal sociology standpoint, I consider this to be an important critical juncture that set countries in a path of political development. The paper version of the chapter (currently \emph{\input{/Users/hectorbahamonde/RU/Dissertation/Papers/IncomeTaxAdoption/status.txt}\unskip}) can be downloaded \href{http://github.com/hbahamonde/IncomeTaxAdoption/raw/master/Bahamonde_IncomeTaxAdoption.pdf}{\texttt{here}}. 


% papers
In addition to the book manuscript, I am currently expanding the findings of a {\bf series of papers} related to vote-buying and vote-selling, using both {\bf observational} and an original {\bf experimental} designs. 



% Vote-Buying in Brazil// explaining the paper
{\bf Vote-buying.} Going forward, I have a {\bf working \href{https://github.com/hbahamonde/Clientelism_paper/raw/master/Bahamonde_Clientelism_Paper.pdf}{paper}} (\emph{\input{/Users/hectorbahamonde/RU/research/Clientelism_paper/status.txt}\unskip}) on vote-buying in Brazil. The paper starts by recognizing that there is not an agreement on whether parties target groups or individuals. In fact, most of the times, scholars assume that group-targeting and individual-targeting are interchangeable. The paper argues that parties make use of simultaneous segmented targeting techniques. Groups are preferable by brokers when party machines need to secure higher levels of electoral support, relying on the economies of scale and spillover effects that these groups provide. However, individuals are better targets when they are more \emph{identifiable}, that is when poor individuals are nested in non-poor contexts or when non-poor individuals are nested in poor groups. The theory also explains why non-poor individuals are also targeted. The paper uses observational data and matching methods and a short case study (Brazil) to confirm the assumptions made in the models.

% vote-selling // experiments
{\bf Vote-Selling.} With the support of a generous grant, I designed two {\bf experiments} in the U.S. out of a series of experiments to be fielded in Latin America for further comparison. I look at the tipping points where U.S. citizens prefer a monetary incentive rather than keeping their right to choose for whom to vote for. My identification strategy takes advantage of a \emph{list experiment} to capture non-biased answers on socially-condemnable/illegal behaviors (like vote-buying). In a separate study, I designed a \emph{conjoint experiment} to identify which of Dahl's democratic dimensions should fail to predict individual propensities of vote-selling. Conjoint experiments allow researchers to directly isolate complex multi-dimensional concepts (such as \emph{support for democracy}) and observe which dimension(s) is/are associated to the outcome of interest (vote-selling). Preliminary results show that when the \emph{liberal} component, particularly the \emph{right to associate} fails, individuals are more likely to sell. In this sense, this project is very innovative as it departs from the common strategy of correlating the standard support for democracy question with vote-buying. My identification strategy is more comprehensive as it \emph{decomposes} ``support for democracy'' in several dimensions which are theoretically and substantively relevant (i.e., Dahl's democracy dimensions). 


{\bf Summary and Future Research} In sum, my book manuscript, job market paper and working papers on vote-buying and experiments on vote-selling, work toward exploring the effects of inequality on political development from a comparative perspective. I use a widely broad methodological perspective, historical comparisons, time series analyses, and experimental and quasi-experimental methodologies. My goal is to use this toolkit to keep asking ``big'' questions that are fundamental to our discipline. My future research will seek to study the connection between elite competition and democratic regimes, exploring the connection between state building and democratic institutions and its relationship with vertical accountability and the role of the middle class on bureaucratic development. Particularly, I am interested in whether dictatorships have different dynamics of state-building relative to democracies. Please check my \href{http://www.hectorbahamonde.com/teaching/}{\texttt{teaching portfolio}} and see how my {\bf research and teaching interests} match.


% ending
{More information, \href{http://www.hectorbahamonde.com/teaching/}{\texttt{syllabi}}, my \href{http://github.com/hbahamonde/Job_Market/raw/master/Bahamonde_Research_Statement.pdf}{\texttt{research}}, \href{http://github.com/hbahamonde/Job_Market/raw/master/Bahamonde_Teaching_Statement.pdf}{\texttt{teaching}} and \href{http://github.com/hbahamonde/Job_Market/raw/master/Bahamonde_Diversity_Statement.pdf}{\texttt{diversity}} statements, as well as other \href{http://www.hectorbahamonde.com/research/}{\texttt{papers}} are available on my website: \href{http://www.hectorbahamonde.com}{\texttt{www.HectorBahamonde.com}}. Thank you for considering my application. I look forward to hearing from you.\unskip}


\vspace{-10cm}


%----------------------------------------------------------------------------------------

\end{newlfm}
\end{document}
